% Data, verifiers, provenance, and the audit design.

\section{Data, verifiers, and audit design}
\label{sec:data}

\subsection{The board}

We use the Open LLM Leaderboard's \citep{openllmleaderboard, fourrier2024openllm} MATH-Hard
\citep{hendrycksmath2021} results: \NumModels{} models, \NumItems{}
Level-5 items across seven subjects, and the generations the leaderboard itself published,
read from its per-model result repositories. Every model has a response for every item, so
the analysis uses \NumRows{} responses with no missing data and no model dropped. A
prespecified rule covered missingness before the data were assembled, and it did not fire.

\subsection{Three verifiers, and what they share}

The first verifier is the leaderboard's original scorer. Its per-item verdicts are published
in the result files, and we did not have to guess its identity: the result files record the
scoring function that ran, and re-executing that function reproduces the published verdict on
\EmReexecItems{} of \EmReexecItems{} responses. It extracts an answer by matching one
sentence pattern and compares symbolically.

The second is Math-Verify, the verifier behind the re-scoring the leaderboard applied to its
MATH results in February 2025 \citep{kydlicek2025fixing}. The re-scoring rewrote the published
scores while leaving the
per-item verdict files untouched, so the leaderboard currently displays scores from one
verifier and publishes item-level verdicts from another. Running Math-Verify ourselves
reproduces the re-scored published aggregates in \MvAggregateMatch{} of
\NumSubjectsRescored{} model-subject cells, with the remainder differing by a single item.

The third extracts the last boxed expression, which is how the benchmark's own reference
solutions mark their answers, and compares it symbolically after the same normalisation the
harness uses.

These verifiers are not independent, and we do not describe them as such.
Table~\ref{tab:shared} lists what they share: a normalisation routine, a symbolic algebra
library, and the same reference answers. What differs is the extraction rule, the parser, and
the source from which the reference answer is taken. Independence is not required for any
result in this paper. What is required is that they fail on different responses, which is an
observable property, and they do.

\subsection{A defect in our own verifier}

While checking provenance we found that the boxed verifier had been silently degraded on this
machine. Its symbolic equivalence check guards itself with a timer available only on Unix. On
Windows the guard raises, the exception is swallowed by a broad handler, and the function
returns ``not equivalent'' for every call. Every verdict it had produced was therefore a
string comparison.

The corrected verifier changes the count of responses on which it disagrees with the
leaderboard's scorer, over the subset of models analysed in earlier work, from
\DisagreeSbEmSeventeenString{} to \DisagreeSbEmSeventeen{}, and the certified count from
\AnchorStringCombined{} to \AnchorSympyCombined{}. All results in this paper use the
corrected verifier. We report the defect because it is the paper's own argument turned on
its author: a verifier that fails silently produces a plausible number, and nothing in the
output says otherwise.

\subsection{Human labels}

Labels come from two audits of the same population. The first is an existing 400-item audit
whose design was registered before labelling; of its labels, \AuditNAgreeAll{} fall in the
cell where every verifier credits the response and \AuditNRejectAll{} where none does. The
second is a \AuditBudget{} label audit designed for this paper and registered before the sheet
was drawn.

Labelling is done in a purpose-built offline tool that shows one response at a time, blocks
pasting, records the time spent on each item, and writes an export whose hash is recorded.
Every batch passes an automated quality gate before it is used. The gate checks structure,
timing, and the rate at which a typed answer exactly matches what a regular expression would
have extracted, compared against the rate observed in the author's earlier audit. Batches that
fail are excluded and retained for the record, and this has happened.

\subsection{Where the labels go}

The design question is not how many labels to buy but where to spend them, and the two
purposes compete. Labels in the discordant cells narrow the identified range and certify more
orderings. Labels in the concordant cells test Assumption~\ref{ass:h}, on which every
certified ordering depends, but do not themselves certify anything.

We computed the trade-off before drawing the sheet, over a grid of allocations, and it is
reported in Figure~\ref{fig:frontier} as a projection rather than as a result. Two findings
came out of it. Certified orderings are nearly flat across allocations, because the binding
constraints are the pooled concordant bound and item sampling rather than discordant coverage;
the whole \AuditBudget{} label budget moves the projected count from \CertPrereg{} to
\ProjectedCertPrereg{}. And a screen powerful enough to detect a moderate departure from
Assumption~\ref{ass:h} consumes most of the budget for a small number of models.

Since certification is nearly insensitive to the split, the budget goes where it buys the most
information about the assumption. The registered design screens the \ScreenedModels{} models
appearing most often in the orderings with the smallest breakdown values, with
\ScreenNRejectAll{} labels in the cell where no verifier credits and \ScreenNAgreeAll{} where
all three do, sized for \ScreenPowerPct\% power against a per-model departure of \ScreenRho{}. The
remaining \LabelsCertification{} labels go to discordant cells in proportion to their size.

What this screen cannot do is stated with it. It does not detect departures smaller than
\ScreenRho{} at that power, and it says nothing at all about the models it does not screen.
The preregistered rejection rule compares each screened model's error lower bound against the
upper bound of the pooled rate computed without that model's own labels, so that a model
cannot inflate the bound it is tested against.
